% ============================================================
% ฉบับภาษาไทย (short paper) จัดรูปแบบตามตัวอย่างภาษาไทยเดิม:
% คอลัมน์เดียว, Angsana New, ตารางที่/รูปที่, เอกสารอ้างอิง
% ชื่อหัวข้อคงเป็นภาษาอังกฤษตามต้นฉบับภาษาอังกฤษ
% คอมไพล์: cd docs/paper/ieee/th && xelatex main_th && bibtex main_th && xelatex main_th && xelatex main_th
% วินัย claim: ทุกประโยคที่เป็นผลอยู่ในขอบเขตของ docs/CLAIMS.md
% ข้อความสีแดง [..] คือส่วนที่รอผลการทดลองเพิ่มเติม
% ============================================================
\documentclass[12pt]{article}

\usepackage{fontspec}
\usepackage[a4paper,margin=1in]{geometry}
\usepackage{xcolor}
\usepackage{booktabs}
\usepackage{graphicx}
\usepackage{amsmath}
\usepackage{caption}
\usepackage{cite}
\usepackage[colorlinks=true,allcolors=blue!50!black]{hyperref}

% Angsana New วาดอักษรเล็กกว่าขนาดที่ระบุ จึงปรับ Scale ให้อ่านได้ราว 16 pt
\setmainfont{Angsana New}[Scale=1.32]
\setmonofont{Consolas}[Scale=MatchLowercase]
\linespread{1.12}
\XeTeXlinebreaklocale "th"
\XeTeXlinebreakskip = 0pt plus 1pt

\newcommand{\todo}[1]{\textcolor{red!70!black}{\textbf{[#1]}}}
\newcommand{\code}[1]{\texttt{#1}}
\graphicspath{{../../figures/}}
\newcommand{\figlang}{th}\newcommand{\figwidth}{4.6in}\newcommand{\fit}[1]{\resizebox{\ifdim\width>\columnwidth \columnwidth\else\width\fi}{!}{#1}}

\renewcommand{\tablename}{ตารางที่}
\renewcommand{\figurename}{รูปที่}
\renewcommand{\refname}{เอกสารอ้างอิง}
\captionsetup{labelsep=space, font=small}
\renewcommand{\abstractname}{บทคัดย่อ}

\makeatletter
\newcommand{\sectionnumpunct}{.}
\renewcommand{\@seccntformat}[1]{%
  \csname the#1\endcsname\csname #1numpunct\endcsname\hspace{0.5em}}
\makeatother
\setlength{\parindent}{1.5em}

\title{\Large\bfseries Absent, Not Misread:\\
Where Thai OCR Vision-Language Models\\ Lose Vowel and Tone Marks}

\author{%
  ธนกฤต เสมอหน้า\textsuperscript{1} \quad ธนวัฒน์ สุริยะ\textsuperscript{2} \quad ผศ.ดร.\ นนท์ คนึงสุขเกษม\textsuperscript{3}\\[4pt]
  \normalsize
  คณะเทคโนโลยีสารสนเทศ สถาบันเทคโนโลยีพระจอมเกล้าเจ้าคุณทหารลาดกระบัง\\
  เลขที่ 1 ซอยฉลองกรุง 1 แขวงลาดกระบัง เขตลาดกระบัง กรุงเทพมหานคร 10520\\[2pt]
  \small
  \textsuperscript{1}\texttt{67070060@kmitl.ac.th} \quad
  \textsuperscript{2}\texttt{67070236@kmitl.ac.th} \quad
  \textsuperscript{3}\texttt{nont@it.kmitl.ac.th}
}
% รายชื่อผู้เขียนและอาจารย์ที่ปรึกษาคัดจากร่างก่อนหน้า (07-09-2026)
\date{}

\begin{document}
\maketitle

\begin{abstract}
\noindent
Vision-Language Model (VLM) แบบเปิดสำหรับ OCR ภาษาไทยอ่านหน้าเอกสารได้ถูกต้องเกือบทั้งหน้าแล้ว แต่ยังมีคำถามที่สำคัญในทางปฏิบัติอยู่ข้อหนึ่ง คือ สระและวรรณยุกต์ของไทยมีขนาดเล็ก และเครื่องหมายเพียงตัวเดียวก็เปลี่ยนความหมายของคำได้ งานนี้จึงถามว่าเครื่องหมายเหล่านี้หายไปที่ใด ไม่ได้ถามเพียงว่าหายไปเท่าไร เราเสนอคะแนนที่วัดเครื่องหมายโดยไม่หักคะแนนเรื่องลำดับการอ่าน หักคะแนนข้อความที่เขียนเกิน และไม่ให้คะแนนกับข้อความที่ซ้ำจากการวนซ้ำ บนชุด calibration ของ ThaiOCRBench โมเดล Typhoon OCR 1.5 (2B) เขียนวรรณยุกต์ผิดเพียง 0.65\% ในตำแหน่งที่อ่านพยัญชนะฐานถูกต้อง และเมื่อให้ข้อความก่อนหน้าที่ถูกต้อง โมเดลเลือกวรรณยุกต์ถูกต้อง 98.8\% แม้ไม่เห็นภาพ เครื่องหมายที่หายไปส่วนใหญ่ไม่ได้ถูกอ่านเลย กล่าวคือ บรรทัดหายไปทั้งบรรทัด ข้อความขาดเป็นช่วง หรือ output ถูกการวนซ้ำกินไป ในบรรดาวิธี decode ที่ไม่ต้อง train เพิ่ม การหยุดเมื่อตรวจพบการวนซ้ำไม่เปลี่ยน output ที่ไม่วนซ้ำเลยและเพิ่มคะแนน mark F1 ขึ้น 0.8 จุดสำหรับ Typhoon และ 10.8 จุดสำหรับโมเดลไทยขนาด 0.9B ส่วน repetition penalty ทำให้ F1 ของหน้าที่ไม่วนซ้ำลดลง และการหักความรู้ภาษาออกทำให้วรรณยุกต์ที่ถูกอยู่แล้วหายไป \todo{ผลของโมเดลอื่นและการยืนยันบนชุดข้อมูลที่เก็บไว้ (locked split) ยังไม่เสร็จ} ผลทั้งหมดเป็นผลเบื้องต้น ใช้ benchmark เดียว ชุดข้อมูล calibration และกฎบางข้อกำหนดหลังจากเห็นข้อมูลแล้ว

\end{abstract}

\noindent\textbf{\textit{คำสำคัญ}}---OCR ภาษาไทย, vision-language model, วรรณยุกต์และสระ, การวนซ้ำ, การ decode โดยไม่ต้อง train เพิ่ม


\section{Introduction}
อักษรไทยเขียนพยัญชนะเป็นแนวหลัก โดยมีสระและวรรณยุกต์วางซ้อนอยู่ด้านบนและด้านล่าง เครื่องหมายเพียงตัวเดียวก็เปลี่ยนความหมายและเสียงของคำได้ เช่น ``ข่าว'' กับ ``ข้าว'' การถอดข้อความที่ทำเครื่องหมายหายหรือสลับตัว จึงผิดสำหรับผู้อ่านแม้ตัวอักษรอื่นเกือบทั้งหมดจะถูกต้อง อย่างไรก็ตาม การให้คะแนน OCR ภาษาไทยในปัจจุบันใช้ตัวชี้วัดของข้อความทั้งก้อน เช่น ความทับซ้อนของ n-gram และ normalized edit distance ใน ThaiOCRBench \cite{nonesung2025thaiocrbench} หรือ character error rate ในงานอื่น ตัวชี้วัดเหล่านี้บอกไม่ได้ว่าปัญหาอยู่ที่เครื่องหมายหรือไม่ และคะแนนที่เทียบทั้งลำดับยังลงโทษโมเดลที่อ่านคอลัมน์ของหน้าในลำดับต่างจากเฉลยด้วย

ปัจจุบันมี vision-language model (VLM) แบบเปิดที่ฝึกมาสำหรับ OCR ภาษาไทยแล้ว \cite{nonesung2026typhoonocr,nectec2026pathumma,pipatanakul2026synthetic} เมื่อโมเดลเหล่านี้ทำเครื่องหมายผิด มีคำอธิบายที่เป็นไปได้สามข้อ หลักฐานในภาพของเครื่องหมายขนาดเล็กอาจอ่อน ความรู้ภาษาของโมเดลอาจทับหลักฐานในภาพ หรือโมเดลอาจไม่ได้อ่านเครื่องหมายนั้นเลย เพราะข้ามบรรทัด หาข้อความไม่เจอ หรือติดการวนซ้ำ แต่ละข้อต้องแก้ต่างกัน ถ้าหลักฐานในภาพอ่อนก็ควรอ่านซ้ำหรือขยายภาพ ถ้าความรู้ภาษาทับหลักฐานก็ควรหักความรู้ภาษาออกแบบ contrastive decoding \cite{leng2024vcd,favero2024m3id} และถ้าข้อความไม่ถูกอ่านก็ต้องแก้ที่วิธี decode หรือการหาตำแหน่งข้อความ

เราศึกษาคำถามนี้บน ThaiOCRBench ด้วย Typhoon OCR 1.5 และ checkpoint แบบเปิดอีกห้าตัว ในงานอ่านทั้งหน้าและงานอ่านข้อความสั้น งานนี้มีส่วนร่วมสามข้อ
\begin{enumerate}
\item คะแนนเครื่องหมายที่ไม่สนลำดับการอ่าน (หัวข้อ~\ref{sec:metric}) ให้คะแนนสระบน สระล่าง และวรรณยุกต์ทีละบรรทัดของเฉลย ไม่ว่าบรรทัดนั้นจะอยู่ตรงไหนของ output ข้อความที่เขียนเกินถูกหักคะแนน และสำเนาที่เกิดจากการวนซ้ำไม่ได้คะแนน
\item การวินิจฉัยว่าเครื่องหมายหายไปที่ใด (หัวข้อ~\ref{sec:where}) ในทั้งหกโมเดล การไม่ได้อ่านบรรทัดทำให้เสีย mark มากกว่าการอ่าน mark ผิดหลังอ่านพยัญชนะถูก 7 ถึง 12 เท่า สำหรับ Typhoon หลังอ่านพยัญชนะถูกจะเขียนวรรณยุกต์ผิดเพียง 0.65\% และเมื่อได้รับข้อความก่อนหน้าที่ถูกต้อง โมเดลเลือกวรรณยุกต์ถูก 98.8\% แม้ไม่เห็นภาพ \todo{การทดสอบเรื่องภาพกับโมเดลอื่น}
\item การประเมินวิธี decode ที่ไม่ต้องฝึกโมเดลเพิ่ม (หัวข้อ~\ref{sec:decoding}) การหยุดเมื่อตรวจพบการวนซ้ำเพิ่ม mark F1 ให้ทุกโมเดล และแทบไม่เปลี่ยน output ที่จบเองโดยไม่วนซ้ำ ส่วน repetition penalty ทำให้ F1 ของหน้าที่ไม่วนซ้ำลดลง และการหักความรู้ภาษาออกทำให้วรรณยุกต์ที่ถูกอยู่แล้วหายไป \todo{การยืนยันบน locked split}
\end{enumerate}
สำหรับโมเดลที่เราศึกษา ความผิดของเครื่องหมายที่ยังเหลือในโมเดล OCR ภาษาไทยที่เก่งอยู่แล้วเกิดขึ้นก่อนจะถึงตัวเครื่องหมาย วิธี decode และการหาข้อความให้เจอจึงสำคัญกว่าการทำให้มองวรรณยุกต์ชัดขึ้น

\section{Related Work}
\textbf{OCR ภาษาไทยด้วย VLM} ThaiOCRBench \cite{nonesung2025thaiocrbench} ประเมิน VLM ทั้งแบบปิดและแบบเปิดบนงานเอกสารภาษาไทย รวมถึงงานอ่านทั้งหน้าและงานอ่านข้อความสั้น โดยใช้ตัวชี้วัดที่อิงความทับซ้อนของ n-gram และ edit distance ได้แก่ BLEU, METEOR, F1 และ normalized Levenshtein distance ส่วนงานอ่านข้อความสั้นใช้ ANLS บทนำของงานนี้ระบุว่าอักขระซ้อนเป็นความยากของภาษาไทย แต่ไม่มีคะแนนที่แยกเครื่องหมายออกมา หลังจากนั้นมี OCR ภาษาไทยแบบเปิดออกมาหลายตัว Typhoon OCR \cite{nonesung2026typhoonocr} เวอร์ชัน 1.5 เป็น full fine-tune ของ Qwen3-VL-2B \cite{bai2025qwen3vl} และผู้พัฒนาระบุว่า infographic ยังต้องปรับปรุง Pathumma-llm-vision 3.0 \cite{nectec2026pathumma} เป็น fine-tune ของ Qwen3-VL-2B อีกตัว ส่วน wayu-paxa-ocr-zero \cite{pipatanakul2026synthetic} fine-tune PaddleOCR-VL \cite{cui2025paddleocrvl} ด้วยหน้าเอกสารไทยสังเคราะห์ ทุกงานรายงานความแม่นยำรวม ส่วนงานของเราถามว่าเครื่องหมายหายไปที่ใด

\textbf{Diacritic ใน OCR} อักษรที่มี diacritic มากเป็นที่ทราบกันว่ายากสำหรับ OCR มีงานเสนอการแก้ข้อความหลัง OCR สำหรับ diacritic ภาษาเวียดนาม \cite{do2025reference} และงานประเมิน OCR ภาษาอาหรับแยกข้อความที่มี diacritic เป็นอีกเงื่อนไขหนึ่ง \cite{elsafty2025arabic} เราไม่พบงานที่ให้คะแนนเครื่องหมายภาษาไทยแยกจากลำดับการอ่านและข้อความที่เขียนเกิน

\textbf{การวนซ้ำและการตกหล่น} การ decode ที่เลือกคำตามความน่าจะเป็นสูงสุดทำให้ข้อความวนซ้ำ \cite{holtzman2020curious} และโมเดล OCR เอกสารก็มีพฤติกรรมเดียวกัน Nougat รายงานหน้าที่โมเดลพิมพ์ประโยคเดิมซ้ำ และหยุด generate เมื่อสัญญาณความแปรปรวนของ logit ลดลง \cite{blecher2024nougat} olmOCR สุ่มใหม่ด้วย temperature ที่สูงขึ้น \cite{poznanski2025olmocr} pipeline ของ PaddleOCR-VL ตัดข้อความซ้ำหลัง decode เสร็จ \cite{paddlex2025truncate} และโค้ด DeepSeek-OCR ที่เผยแพร่ห้ามการซ้ำของ n-gram ภายในหน้าต่างหนึ่ง \cite{deepseekocr2025code,wei2025deepseekocr} CC-OCR นับการวนซ้ำเป็น hallucination แบบหนึ่งของ OCR \cite{yang2025ccocr} งานเหล่านี้มองการวนซ้ำเป็นสิ่งที่ต้องตรวจจับหรือหลีกเลี่ยง ส่วนเราวัดว่าวิธีแก้แต่ละแบบทำให้เครื่องหมายในหน้าที่ไม่วนซ้ำเสียไปเท่าใด หนึ่งในวิธีที่วัดคือ repetition penalty ของ CTRL \cite{keskar2019ctrl} ซึ่ง package OCR บางตัวเปิดไว้เป็นค่าเริ่มต้น

\textbf{ความรู้ภาษาและ decoding ที่ไม่ต้องฝึกเพิ่ม} Contrastive decoding ลด hallucination ของ VLM โดยลดน้ำหนักสิ่งที่โมเดลจะทำนายเมื่อไม่มีภาพหรือเมื่อภาพถูกทำให้เสื่อม เช่น VCD \cite{leng2024vcd}, M3ID \cite{favero2024m3id} และ PAI \cite{liu2024pai} วิธีเหล่านี้สมมติว่าความรู้ภาษาเป็นต้นเหตุ แต่สำหรับวรรณยุกต์ไทยเราพบว่าความรู้ภาษามักถูกอยู่แล้ว วิธีที่เราใช้เป็นการเขียนใหม่ จึงเรียกว่า M3ID-form และ PAI-style ส่วน speculative decoding \cite{leviathan2023speculative} ที่ใช้ prompt lookup เป็นตัวร่าง \cite{saxena2023pld} ทำให้ decode เร็วขึ้น และ output ไม่เปลี่ยนเมื่อคำนวณแบบแม่นยำ

\section{Scoring Marks Without Charging Reading Order}\label{sec:metric}
เครื่องหมายที่เราวัดเรียกว่า mark ได้แก่ สระบนหกตัว (\textit{กั กิ กี กึ กื ก็}) สระล่างสองตัว (\textit{กุ กู}) และวรรณยุกต์สี่ตัว (\textit{ก่ ก้ ก๊ ก๋}) เฉลยของหน้าเต็มมีหลายบรรทัด และโมเดลอาจอ่านคอลัมน์ของหน้าในลำดับต่างจากเฉลย การจัดแนวแบบ global จึงนับ mark ที่อ่านถูกว่าเป็นข้อผิด

\textbf{นิยาม} เราพิจารณาบรรทัดของเฉลยที่ยาวอย่างน้อยแปดตัวอักษร เรียงจากยาวไปสั้น แต่ละบรรทัดหาช่วงของ output ที่ตรงกับบรรทัดนั้นที่ character error rate ต่ำกว่า 40\% โดยไม่นับอักขระที่ถูกจองไปแล้ว จากนั้นให้คะแนน mark ที่ตรงกัน และจองช่วงนั้นไว้ เพื่อให้ช่วงเดียวของ output ไม่ให้คะแนนสองบรรทัด บรรทัดที่เหลือถูกจัดแนวแบบ global ตามลำดับของเฉลยกับ output ที่ตัดอักขระที่จองไว้ออกทั้งหมด หากไม่มีขั้นตอนที่สองนี้ mark ในบรรทัดที่อ่านคลาดเคลื่อนจะหลุดไป และเวอร์ชันแรกของเราไม่ผ่านการตรวจกับโมเดลตั้งต้น Recall คือ mark ที่ได้คะแนนหารด้วย mark ในเฉลย ส่วน precision คือ mark ที่ได้คะแนนหารด้วย mark ทั้งหมดใน output ซึ่งทำให้ข้อความที่เขียนเกินถูกหักคะแนน ส่วน F1 คำนวณจากผลรวมของทุกข้อ

\textbf{แบบที่กันการวนซ้ำ} ในแบบนี้ คะแนนคำนวณหลังย่อทุกช่วงที่วนซ้ำให้เหลือหนึ่งสำเนา ช่วงที่วนซ้ำคือหน่วยยาว 4 ถึง 400 ตัวอักษรที่มีตัวอักษรอย่างน้อยหนึ่งตัวและซ้ำกันเหมือนกันทุกตัวอย่างน้อยแปดครั้ง หรือหกครั้งสำหรับหน่วยที่ยาว 50 ตัวอักษรขึ้นไป ส่วน precision ยังนับ mark ทุกตัวของ output เต็ม สำเนาจึงถูกหักคะแนนแต่ไม่ได้ recall เราเพิ่มแบบนี้หลังจากผู้ร่วมงานพบว่า alignment จับคู่สำเนาของหน่วยเดียวที่ซ้ำ 434 ครั้งเข้ากับพยางค์ทั่วไปในเฉลย ถ้าไม่มีการวนซ้ำ คะแนนทั้งสองแบบเท่ากัน

\textbf{การตรวจสอบ} การอ่านหน้าเดียวกันสองรอบต่อกันได้คะแนนต่ำกว่าการอ่านรอบเดียวที่ดีกว่ามาก recall ของโมเดลตั้งต้นต่างจาก global alignment ไม่เกิน 0.3 จุดในทุกเซลล์ และ recall ของ Typhoon บนหน้าเต็มเพิ่มขึ้น 1.6 ถึง 1.8 จุดเมื่อไม่หักคะแนนเรื่องลำดับ การหักเรื่องลำดับอย่างเดียวทำให้ mark F1 ของ Typhoon บนหน้าเต็มลดจาก 94.9 เหลือ 93.1 สำหรับ prompt BQ (จาก 95.3 เหลือ 93.7 สำหรับ TC ซึ่งนิยามในหัวข้อถัดไป) โมเดลตั้งต้นได้รับผลกระทบน้อยมาก

\section{Experimental Setup}\label{sec:setup}
\textbf{ข้อมูล} เราใช้ชุด test ของ ThaiOCRBench \cite{nonesung2025thaiocrbench} สองงาน คืองานอ่านทั้งหน้าและงานอ่านข้อความสั้น ชุด calibration สุ่มแบบกำหนด seed ไว้ 30\% แบ่งชั้นตามงานและหมวด มีงานอ่านทั้งหน้า 69 ข้อ และงานอ่านข้อความสั้น 109 ข้อ ส่วนที่เหลืออีก 70\% เป็นชุด locked ซึ่งยังไม่ถูกเปิด เรารายงาน output และคะแนนรายข้อ และไม่เผยแพร่ภาพของ benchmark

\textbf{โมเดล} เราประเมิน Typhoon OCR 1.5 (2B) โมเดลตั้งต้น Qwen3-VL-2B, Qwen3-VL-4B และโมเดลขนาด 0.9B ที่ออกแบบให้อ่านทีละกรอบข้อความ คือ PaddleOCR-VL-1.6 และ wayu-paxa-ocr-zero (เรียกว่า Wayu) โดยนำสองตัวหลังมาอ่านทั้งหน้าด้วย ซึ่งเกินกว่าที่ออกแบบไว้ ทุก checkpoint ระบุ revision ไว้ และเราใช้ revision ของ Typhoon ที่เผยแพร่ภายใต้ Apache-2.0 อย่างเดียว ค่า revision, seed และรหัสการรันบันทึกไว้ใน repository ของโครงการ

\textbf{Prompt และการ decode} BQ คือคำถามของ benchmark ในแต่ละข้อ ส่วน TC คือ prompt ที่ model card ของ Typhoon แนะนำ Paddle และ Wayu ใช้ prompt \code{OCR:} การ decode ใช้ greedy สร้าง token ใหม่ไม่เกิน 3,072 ตัว ด้วย fp16 บน GPU Kaggle T4 แต่ละผลมาจากการรันครั้งเดียว output ถูกแปลงเป็นข้อความธรรมดาก่อนให้คะแนน BQ เป็น prompt หลัก และ TC เป็น prompt รอง

\textbf{สถิติ} ช่วงความเชื่อมั่นคือ bootstrap 95\% ระดับรายข้อ ผลทั้งหมดเป็นผลเบื้องต้นจากชุด calibration และกฎ decode บางข้อถูกกำหนดหลังจากเห็นชุดนี้แล้ว เราระบุไว้ในจุดที่เกี่ยวข้อง

\section{Where Marks Are Lost}\label{sec:where}
\subsection{Six Models}
Typhoon ได้ mark F1 94.9 บนหน้าเต็ม และ 81.7 บนงานอ่านข้อความสั้น (รูปที่~\ref{fig:models} และตารางที่~\ref{tab:models}) Pathumma-llm-vision 3.0 \cite{nectec2026pathumma} ซึ่งเป็น fine-tune ภาษาไทยตัวที่สองจากโมเดลตั้งต้นเดียวกัน ได้ใกล้เคียงกัน เมื่อใช้ prompt ของ model card เอง ได้ 89.8 และ 77.6 ส่วนต่างจาก Typhoon คือ $-$5.1 จุดบนหน้าเต็ม (ช่วง $-$14.3 ถึง +1.6) และ $-$4.1 บนงานอ่านข้อความสั้น (ช่วง $-$15.1 ถึง +8.0) เราจึงบอกไม่ได้ว่าตัวใดดีกว่า ทั้งสองตัวสูงกว่าโมเดลตั้งต้นมาก คือ 48.9 และ 53.3 จุด ส่วน checkpoint อีกสามตัวต่ำกว่า Typhoon มาก เช่น Wayu ต่ำกว่า 29.6 จุดบนหน้าเต็ม (ช่วง 19.2 ถึง 39.9) การเพิ่มขนาดโมเดลตั้งต้นเป็นสองเท่าเปลี่ยน F1 บนหน้าเต็มเพียง +5.0 (ช่วง $-$4.0 ถึง +14.1) และอัตราผิดของวรรณยุกต์ยังอยู่ราว 9.4 ถึง 9.9\% ค่ามัธยฐานของ character error rate ต่อหน้าของ Typhoon คือ 2.8\% โมเดลทั้งหมดถูกเทียบตามที่เผยแพร่ ซึ่งจำนวน visual token ไม่เท่ากัน Typhoon แข็งที่สุดในบรรดาโมเดลที่เราวัด และใกล้เคียงกับ Pathumma model card ของ Pathumma รายงานคะแนน ThaiOCRBench และไม่ระบุที่มาของข้อมูลฝึก 377K ตัวอย่าง เราจึงตัดความเป็นไปได้ที่ข้อมูลซ้อนกับ benchmark ไม่ได้

ในงานอ่านข้อความสั้น สองโมเดลพลาดต่างกัน Pathumma พบ mark ที่ขอได้มากกว่า (recall สูงกว่า 14.3 จุด) แต่มักถอดทั้งภาพแทนที่จะถอดเฉพาะช่วงที่ถาม ทำให้ precision ต่ำลง

\begin{figure}[t]
\centering\includegraphics[width=\figwidth]{fig_models_\figlang.pdf}
\caption{Mark F1 บนชุด calibration (ก) งานอ่านทั้งหน้า (ข) งานอ่านข้อความสั้น ใช้ prompt BQ ยกเว้นโมเดล 0.9B สองตัวที่ใช้ \code{OCR:} และ Pathumma ที่ใช้ prompt ของ model card}\label{fig:models}
\end{figure}

\begin{table}[t]
\caption{Mark F1 อัตราผิดของวรรณยุกต์หลังอ่านพยัญชนะถูก และสัดส่วน output ที่ชนขีดจำกัด token (\%) Pathumma ใช้ prompt ของ model card}\label{tab:models}
\centering\footnotesize
\fit{\begin{tabular}{@{}l rrr rr@{}}
\toprule
 & \multicolumn{3}{c}{อ่านทั้งหน้า} & \multicolumn{2}{c}{อ่านข้อความสั้น}\\
\cmidrule(lr){2-4}\cmidrule(l){5-6}
โมเดล & F1 & ผิดวรรณยุกต์ & ชนขีดจำกัด & F1 & ชนขีดจำกัด\\
\midrule
Typhoon OCR 1.5 & \textbf{94.9} & 0.65 & 2.9 & \textbf{81.7} & 0.9\\
Pathumma 3.0 & 89.8 & 0.99 & 5.8 & 77.6 & 1.8\\
Wayu & 65.2 & 5.27 & 26.1 & 46.0 & 20.2\\
PaddleOCR-VL-1.6 & 53.3 & 14.57 & 39.1 & 47.9 & 14.7\\
Qwen3-VL-4B & 45.9 & 9.89 & 27.5 & 23.5 & 11.9\\
Qwen3-VL-2B & 41.0 & 9.41 & 43.5 & 24.3 & 14.7\\
\bottomrule
\end{tabular}}
\end{table}

\textbf{สำเนาของการวนซ้ำเคยได้ recall} คะแนนแบบกันการวนซ้ำเปลี่ยน Typhoon ไม่เกิน 0.2 จุด แต่ทำให้ F1 บนหน้าเต็ม (BQ) ของโมเดลตั้งต้นลดจาก 41.0 เหลือ 34.4 และ recall ลดจาก 65.0 เหลือ 54.6 (ตารางที่~\ref{tab:repeat}) ราวหนึ่งในหกของ mark ที่คะแนนแบบไม่สนลำดับให้แก่โมเดลตั้งต้นมาจากสำเนาของการวนซ้ำ เราเปิดดู output ของทุกช่วงที่ถูกตัดแล้ว ทั้ง 126 ช่วงชนขีดจำกัด token ช่วงที่สั้นที่สุดมี 15 สำเนา และไม่มีข้อความปกติถูกตัดออก

\begin{table}[t]
\caption{Recall (R), precision (P) และ F1 (\%) เมื่อให้และไม่ให้คะแนนแก่สำเนาของการวนซ้ำ หน้าเต็ม BQ}\label{tab:repeat}
\centering\footnotesize
\fit{\begin{tabular}{@{}l rrr rrr@{}}
\toprule
 & \multicolumn{3}{c}{ไม่สนลำดับ} & \multicolumn{3}{c}{กันการวนซ้ำ}\\
\cmidrule(lr){2-4}\cmidrule(l){5-7}
โมเดล & R & P & F1 & R & P & F1\\
\midrule
Typhoon & 96.1 & 93.7 & 94.9 & 96.1 & 93.6 & 94.8\\
Qwen3-VL-2B & 65.0 & 29.9 & 41.0 & 54.6 & 25.1 & 34.4\\
\bottomrule
\end{tabular}}
\end{table}

\subsection{Most Lost Marks Are Never Read}
บนหน้าเต็ม Typhoon ทำ mark ผิด 6.2\% (BQ) และ 7.1\% (TC) ของ mark ในเฉลยทั้งหมด เราจัดแต่ละข้อผิดเข้าสาเหตุเดียวโดยอัตโนมัติ (รูปที่~\ref{fig:causes}) บรรทัดอาจถูกอ่านในลำดับอื่น คือมีช่วงของ output ตรงกับบรรทัดนั้นที่ character error rate ต่ำกว่า 20\% (สำหรับบรรทัดที่ยาวอย่างน้อยแปดตัวอักษร) บรรทัดอาจหายไปทั้งบรรทัด ข้อความอาจขาดเป็นช่วงภายในบรรทัดที่เก็บไว้ หรือ mark อาจถูกอ่านผิดตัว บรรทัดที่หายทั้งบรรทัดคิดเป็น 2.3\% (BQ) และ 2.9\% (TC) ของ mark ข้อความที่ขาดเป็นช่วงคิดเป็น 0.7\% และ 0.8\% และการอ่านผิดตัวคิดเป็น 1.5\% และ 2.4\% ส่วนลำดับการอ่านอย่างเดียวคิดเป็น 1.6\% และ 0.9\% ซึ่งไม่ใช่ความผิดพลาดของการอ่าน เมื่อเปิดดู output จริง บรรทัดที่หายส่วนใหญ่เป็นข้อความใน infographic และกราฟ โฆษณาที่แทรกอยู่ รวมถึงชื่อเรื่องและส่วนประกอบของหน้า ซึ่ง Typhoon เขียนเป็นตัวแทนภาพไว้ สอดคล้องกับที่ผู้พัฒนาระบุเรื่อง infographic \cite{nonesung2026typhoonocr} การจัดสาเหตุเป็นแบบอัตโนมัติ และเราตรวจภาพหน้าเอกสารเฉพาะหกหน้าที่มี mark ที่หายเป็นส่วนใหญ่

รูปแบบเดียวกันพบในอีกห้าโมเดล (รูปที่~\ref{fig:causes}) เราจัดสาเหตุบน output ที่ตัดที่การวนซ้ำครั้งแรก เพราะ output ที่วนซ้ำยาวกว่าเฉลยหลายเท่า และการจัดแนวเชื่อถือไม่ได้ บรรทัดที่หายทั้งบรรทัดรวมกับข้อความที่ขาดเป็นช่วงทำให้เสีย mark 3.2\% สำหรับ Typhoon 9.3\% สำหรับ Pathumma และ 29 ถึง 44\% สำหรับอีกสี่โมเดล ส่วนการอ่าน mark ผิดหลังอ่านพยัญชนะถูกทำให้เสีย 0.4\% สำหรับ Typhoon 0.8\% สำหรับ Pathumma และ 2.7 ถึง 4.3\% สำหรับโมเดลที่เหลือ ในทุกโมเดล การไม่ได้อ่านจึงทำให้เสีย mark มากกว่าการอ่านผิด 7 ถึง 12 เท่า สำหรับโมเดลที่วนซ้ำ บรรทัดที่อยู่หลังการวนซ้ำถูกนับว่าหายไป และเราแยกไม่ได้ว่าบรรทัดเหล่านั้นหายเพราะการวนซ้ำหรือเพราะโมเดลไม่อ่านอยู่แล้ว

\begin{figure}[t]
\centering\includegraphics[width=\figwidth]{fig_causes_\figlang.pdf}
\caption{mark บนหน้าเต็มที่หายไปกับสาเหตุใด คิดเป็นสัดส่วนของ mark ในเฉลยทั้งหมด สำหรับ output ที่ตัดที่การวนซ้ำครั้งแรก ยอดรวมของแต่ละแท่งไม่รวมส่วนสีเทา ซึ่งเป็นเพียงความต่างของลำดับการอ่านและไม่ใช่ข้อผิด}\label{fig:causes}
\end{figure}

\subsection{With the Right Context, the Prior Is Already Right}
ในการทดสอบแบบเลือกคำตอบบังคับ เราป้อนข้อความที่ถูกต้องจนถึงตำแหน่งของ mark แล้วให้โมเดลให้คะแนนตัวเลือกที่เป็นไปได้ของ mark นั้นทั้งแบบมีภาพและไม่มีภาพ (fp32, tokenize ในบริบทเดียวกัน) บนหน้าเต็ม Typhoon เลือกวรรณยุกต์ถูก 99.7\% เมื่อมีภาพ และ 98.8\% เมื่อไม่มีภาพ ส่วนโมเดลตั้งต้นได้ 98.1\% และ 98.9\% ภาพจึงช่วยเพิ่ม 0.9 จุดให้ Typhoon และลดลง 0.8 จุดสำหรับโมเดลตั้งต้น (รูปที่~\ref{fig:prior}ก) การทดสอบนี้ให้บริบทที่ถูกต้องแก่โมเดล จึงไม่ใช่การ generate จริง ในคู่โมเดลนี้เพียงคู่เดียว โมเดลเฉพาะทางได้ประโยชน์จากภาพมากกว่าโมเดลตั้งต้น และเราไม่ขยายความไปเกินกว่านี้ \todo{การทดสอบเดียวกันกับ Pathumma และ Qwen3-VL-4B}

\textbf{ความมั่นใจของโมเดล} ความน่าจะเป็นของ token ที่ Typhoon เขียนเองบอกตำแหน่งที่อ่าน mark ผิดได้ ค่าความน่าจะเป็นต่ำสุดของ token ภายในกลุ่มอักขระแยก mark ที่อ่านผิดออกจาก mark ที่ถูกด้วย AUROC 0.918 (BQ) และ 0.935 (TC) บน output แบบ greedy ของโมเดลเอง

\begin{figure}[t]
\centering\includegraphics[width=\figwidth]{fig_prior_\figlang.pdf}
\caption{(ก) ความแม่นยำของวรรณยุกต์ที่เพิ่มขึ้นเมื่อมีภาพ ที่ตำแหน่งซึ่งลงทะเบียนไว้ (ข) ความแม่นยำของ Typhoon เมื่อหักคะแนนแบบไม่มีภาพด้วยน้ำหนัก $\lambda$}\label{fig:prior}
\end{figure}

\section{Training-Free Decoding and Marks}\label{sec:decoding}
\subsection{Loops Are Rare but Costly}
จาก output ของ Typhoon 356 รายการ (ทั้งสอง prompt ทั้งสองงาน) มี 11 รายการที่วนซ้ำ ในแต่ละเซลล์ หน้าที่วนซ้ำเพียงสองถึงห้าหน้ามี mark ที่เขียนเกินรวมกัน 68 ถึง 76\% ของทั้งหมด ถ้าไม่มีหน้าเหล่านี้ F1 บนหน้าเต็ม (BQ) จะเป็น 97.4 แทน 94.9 ตัวเลขนี้เป็นเพดานที่ได้จากการตัดหน้าเหล่านั้นออก ไม่ใช่ผลของวิธีใด

\subsection{Three Ways to Handle a Loop}
เราเทียบวิธีจัดการการวนซ้ำสามวิธีสำหรับ Typhoon กับ greedy decoding (รูปที่~\ref{fig:decoders} หน้าเต็ม BQ)
\begin{itemize}
\item Repetition penalty 1.1 ซึ่งเป็นค่าต้านการวนซ้ำใน package ของ Typhoon เอง ใช้ penalty แบบเดียวกับ CTRL \cite{keskar2019ctrl}
\item n-gram blocking แบบเดียวกับโค้ดของ DeepSeek-OCR ซึ่งห้ามการซ้ำของ 30-gram ภายใน 90 token
\item หยุดเมื่อตรวจพบการวนซ้ำ เมื่อหน่วยข้อความซ้ำแปดครั้ง หรือหกครั้งสำหรับหน่วยที่ยาว 50 ตัวอักษรขึ้นไป การ decode จะหยุดที่สำเนาแรก ผลของ greedy decoding ถูกกำหนดจาก prefix อยู่แล้ว เราจึงคำนวณจุดตัดแบบ offline จาก output greedy ที่ครบแล้ว ส่วนการทำตอน decode จริงยังไม่ได้ตรวจเทียบกับผลนี้
\end{itemize}

\begin{figure}[t]
\centering\includegraphics[width=\figwidth]{fig_decoders_\figlang.pdf}
\caption{การเปลี่ยนแปลงของ mark F1 บนหน้าเต็มเทียบกับ greedy decoding (Typhoon, BQ) แท่งสีน้ำเงินคือทุกหน้า พร้อมช่วงความเชื่อมั่น 95\% แท่งสีส้มคือเฉพาะหน้าที่ไม่วนซ้ำ}\label{fig:decoders}
\end{figure}

\textbf{Repetition penalty} เปลี่ยน output เกือบทุกรายการ หน้าเต็ม BQ ที่ไม่วนซ้ำไม่มีหน้าใดเหมือน greedy เลย และ F1 ของหน้าที่ไม่วนซ้ำลดลง 1.1 ถึง 1.3 จุด ในงานอ่านข้อความสั้น (BQ) อัตราผิดของวรรณยุกต์เพิ่มจาก 17.1\% เป็น 22.6\% เมื่อรวมทุกหน้า ผลแยกจากศูนย์ไม่ได้ (+0.89 ช่วง $-$2.70 ถึง +4.81) ข้อความไทยมีพยางค์และวรรณยุกต์ซ้ำกันอยู่ตามปกติ ซึ่ง penalty ทั้งหน้าอาจกดไว้ แต่เรายังไม่ได้ทดสอบคำอธิบายนี้

\textbf{n-gram blocking} ปลอดภัยเมื่อไม่มีการวนซ้ำ output ที่ไม่วนซ้ำไม่เปลี่ยน 92 ถึง 100\% และ F1 เปลี่ยนไม่เกิน 0.05 วิธีนี้หยุดการวนซ้ำได้ 6 จาก 11 รายการ ในรายการที่เหลือ โมเดลหนีไปวนแบบใกล้เคียงของเดิม เช่น ชื่อซอยที่เพี้ยนไปเป็นชื่อซอยคล้ายกันสามชื่อ ข้อความเหล่านี้มี mark เพิ่ม ทำให้ mark ที่เขียนเกินสูงขึ้น และ F1 บนหน้าเต็มต่ำลง 0.68 สำหรับ BQ เราตีความว่าการวนซ้ำเป็นสถานะดึงดูดของ decoder การห้ามสำเนาที่ซ้ำตรงตัวจึงย้ายโมเดลไปอยู่ในสถานะข้างเคียง นี่เป็นการตีความ ยังไม่ได้ทดสอบ

\textbf{หยุดเมื่อพบการวนซ้ำ} ไม่เปลี่ยน output ที่ไม่วนซ้ำเลยในทุกเซลล์ mark F1 บนหน้าเต็มเพิ่มขึ้น 0.84 (BQ) และ 0.67 (TC) และ F1 ของงานอ่านข้อความสั้นเพิ่มขึ้น 3.30 และ 2.39 การหยุดยังประหยัด token ที่สร้างได้ 11.2\% ส่วน recall ลดลงไม่เกิน 0.16 จุด เพราะการหยุดกู้ข้อความที่การวนซ้ำทำให้ข้ามไปไม่ได้ เราเขียนกฎนี้หลังจากเห็นผลการเปรียบเทียบก่อนหน้า และแต่ละเซลล์มี output ที่วนซ้ำเพียงหนึ่งถึงสี่รายการ ช่วงความเชื่อมั่นทุกช่วงจึงรวมศูนย์ไว้

\subsection{The Stop on Other Models}
โมเดลอื่นวนซ้ำบ่อยกว่า Typhoon มาก บนหน้าเต็ม output ของโมเดลตั้งต้น Qwen3-VL-4B Paddle และ Wayu วนซ้ำ 18 ถึง 30 จาก 69 รายการ เทียบกับ 2 รายการของ Typhoon และ 4 รายการของ Pathumma เมื่อใช้การหยุดแบบ offline เดียวกันกับ output greedy ของโมเดลเหล่านี้ mark F1 เพิ่มขึ้นในทุกโมเดลทั้งสองงาน (ตารางที่~\ref{tab:stop}) ไม่มีช่วงความเชื่อมั่นใดต่ำกว่าศูนย์ และเจ็ดจากสิบสองเซลล์ไม่รวมศูนย์ ผลที่เพิ่มมากในโมเดลที่วนซ้ำบ่อย เช่น +24.4 สำหรับโมเดลตั้งต้น และ +10.8 สำหรับ Wayu บนหน้าเต็ม (ช่วง 5.3 ถึง 16.1) และน้อยสำหรับ Typhoon และ Pathumma ซึ่งวนซ้ำไม่บ่อย ผลที่เพิ่มมาจาก precision และแลกกับ recall ที่ลดลงสูงสุด 10.6 จุดบนหน้าเต็ม recall ส่วนหนึ่งเป็นคะแนนที่สำเนาของการวนซ้ำเคยได้ และที่เหลือคือข้อความหลังการวนซ้ำที่ถูกตัดทิ้ง ในสิบสองเซลล์ การหยุดเปลี่ยน output ที่จบเองเพียง 3 จาก 898 รายการ สองรายการเป็น output ของ Paddle ที่ออกจากการวนซ้ำเองแล้ว และการตัดทำให้ข้อความจริงของหน้าที่ตามหลังหายไป repetition penalty 1.05 เพียงอย่างเดียว ซึ่งเป็นค่าในตัวอย่างของ Wayu เอง ไม่ช่วย (+2.8 ช่วง $-$3.6 ถึง +9.6 บนหน้าเต็ม และ $-$11.3 ช่วง $-$22.5 ถึง +4.5 บนงานอ่านข้อความสั้น) การหยุดประหยัด token ที่สร้างได้ 34\% บนหน้าเต็มของ Wayu ในการรันนี้ 72\% ของช่องว่าง recall ระหว่าง Wayu กับ Typhoon บนหน้าเต็มอยู่ใน 18 หน้าที่ Wayu วนซ้ำ และไม่มีวิธี decode ใดที่เราทดสอบกู้คืนได้ เรายังบอกไม่ได้ว่าการวนซ้ำเป็นสาเหตุหรือเป็นอาการ

\begin{table}[t]
\caption{Mark F1 ก่อนและหลังหยุดเมื่อตรวจพบการวนซ้ำ (แบบ offline)}\label{tab:stop}
\centering\footnotesize
\fit{\begin{tabular}{@{}l rr rr@{}}
\toprule
 & \multicolumn{2}{c}{อ่านทั้งหน้า} & \multicolumn{2}{c}{อ่านข้อความสั้น}\\
\cmidrule(lr){2-3}\cmidrule(l){4-5}
โมเดล & greedy & + หยุด & greedy & + หยุด\\
\midrule
Typhoon OCR 1.5 & 94.9 & 95.7 & 81.7 & 85.0\\
Pathumma 3.0 & 89.8 & 92.8 & 77.6 & 77.7\\
Wayu & 65.2 & 76.0 & 46.0 & 70.4\\
PaddleOCR-VL-1.6 & 53.3 & 58.4 & 47.9 & 61.8\\
Qwen3-VL-4B & 45.9 & 68.8 & 23.5 & 29.9\\
Qwen3-VL-2B & 41.0 & 65.4 & 24.3 & 66.9\\
\bottomrule
\end{tabular}}
\end{table}

\subsection{Subtracting the Language Prior Removes Correct Marks}
ในการทดสอบแบบเลือกคำตอบบังคับ การหักคะแนนแบบไม่มีภาพด้วยน้ำหนัก $\lambda$ ทำให้ความแม่นยำของวรรณยุกต์ของ Typhoon ลดจาก 99.7\% เป็น 97.3\% ที่ $\lambda=0.5$ และเหลือ 42.3\% ที่ $\lambda=1$ (รูปที่~\ref{fig:prior}ข) ส่วนโมเดลตั้งต้นลดเหลือ 84.7\% ที่ 0.5 ในการ generate จริง บนกลุ่มย่อย 24 ข้อ contrast แบบ M3ID-form ของเราหยุดการวนซ้ำของโมเดลตั้งต้นได้ (character error rate จาก 1.12 เหลือ 0.21) แต่ก็ทำให้อัตราผิดของวรรณยุกต์ในข้อที่เก็บไว้เพิ่มจาก 3.7\% เป็น 7.6\% ส่วนใหญ่เพราะวรรณยุกต์ถูกลบ (21 เป็น 55 ใน 12 ข้อ) ส่วน Typhoon เปลี่ยนเพียงเล็กน้อยและกระจุกอยู่ในไม่กี่ข้อ ผลทั้งสองชี้ไปทางเดียวกัน ความรู้ภาษาเกี่ยวกับวรรณยุกต์ถูกอยู่แล้ว การหักออกจึงลบวรรณยุกต์ที่ถูกไปด้วย

\subsection{Speed}
Speculative decoding แบบ prompt lookup ทำให้ Typhoon เร็วขึ้น 1.14 ถึง 1.15 เท่า (โมเดลตั้งต้น 1.16 ถึง 1.21 เท่า) บนหน้าที่ไม่เสื่อมเป็นการวนซ้ำ output เหมือนเดิมทุกตัวอักษรยกเว้นจุดที่คะแนนของสองตัวเลือกเกือบเท่ากันใน fp16 โดย 88 ถึง 93\% ของทุกข้อเหมือนเดิม ตัวอ้างอิง greedy ให้ output ตรงกับ output หลักของเราทุกตัวอักษรในทั้ง 177 ข้อที่จับเวลา \todo{ความเร็วเมื่อใช้การหยุดร่วมกับ speculative decoding}

\section{Discussion and Limitations}
ในทั้งหกโมเดลที่เราศึกษา ความผิดของ mark ส่วนใหญ่เกิดก่อนจะถึงตัว mark ได้แก่ บรรทัดที่ไม่ได้อ่าน ข้อความที่ขาดเป็นช่วง และการวนซ้ำ เมื่อได้บริบทที่ถูกต้อง วรรณยุกต์แทบแน่นอนอยู่แล้ว และวิธีที่หักความรู้ภาษาออกต้องแลกด้วยวรรณยุกต์ วิธีที่ถูกและปลอดภัยในการจัดการการวนซ้ำคือหยุดมัน

งานนี้มีข้อจำกัดหลายข้อ ใช้ benchmark เดียวและสองงาน ทุกผลมาจากชุด calibration และกฎการหยุด คะแนนแบบกันการวนซ้ำ และการให้คะแนนถูกเขียนหลังจากเห็นข้อมูลแล้ว ชุด locked ยังไม่ถูกเปิด \todo{รอดำเนินการ} โมเดลถูกเทียบตามที่เผยแพร่ ซึ่งจำนวน visual token ไม่เท่ากัน และโมเดลที่ออกแบบให้อ่านทีละกรอบสองตัวถูกนำมาอ่านทั้งหน้าเกินกว่าที่ออกแบบไว้ Pathumma อาจเคยเห็น benchmark ตอนฝึก output ที่วนซ้ำของ Typhoon และ Pathumma มีน้อย ช่วงความเชื่อมั่นส่วนใหญ่จึงรวมศูนย์ไว้ ข้อสรุปของเราจึงเป็นเรื่องกลไก ไม่ใช่อัตรา การหยุดบางครั้งตัดข้อความจริงที่ตามหลังการวนซ้ำ และการจัดสาเหตุหลังการตัดแยกการวนซ้ำออกจากการไม่อ่านไม่ได้ การจัดสาเหตุของข้อผิดเป็นแบบอัตโนมัติ การทดสอบแบบเลือกคำตอบบังคับใช้บริบทที่ถูกต้อง และ contrast ที่เราใช้เป็นการเขียนใหม่ที่ทดสอบบน 24 ข้อ สุดท้าย เราใช้ revision ของ Typhoon ที่เผยแพร่ภายใต้ Apache-2.0 อย่างเดียว ขณะที่เงื่อนไขปัจจุบันของผู้พัฒนาขอให้ได้รับความยินยอมเป็นลายลักษณ์อักษรก่อนนำไปเปรียบเทียบประสิทธิภาพ

\section{Conclusion}
OCR ภาษาไทยแบบเปิดที่เก่งที่สุดไม่ได้พลาดที่วรรณยุกต์เป็นหลัก แต่พลาดก่อนหน้านั้น คืออ่านไม่ครบ หาข้อความไม่เจอ หรือวนซ้ำ การหยุดเมื่อตรวจพบการวนซ้ำเป็นขั้นตอนที่ปลอดภัย ไม่ต้องฝึกโมเดลเพิ่ม และให้ผลที่วัดได้ ส่วนการหักความรู้ภาษาออกไม่ใช่ทางที่ดี \todo{ข้อสรุปสุดท้ายรอผลการทดลองที่เหลือ}

\section*{Ethical and Licensing Notes}
โมเดลทั้งหมดเป็นโมเดลแบบเปิดและใช้ตามสัญญาอนุญาตของแต่ละตัว (Apache-2.0 หรือ MIT โดย Typhoon ใช้ revision ที่เป็น Apache-2.0 อย่างเดียว) ThaiOCRBench เผยแพร่ภายใต้ CC-BY-SA-4.0 เราอ้างอิงบทความและ revision ของ benchmark รายงาน output และคะแนน และไม่เผยแพร่ภาพของ benchmark ซ้ำ


\bibliographystyle{IEEEtran}
\bibliography{../../references}

\end{document}
